\documentclass[conference]{IEEEtran}
\usepackage{booktabs}
\usepackage{graphicx}
\usepackage{hyperref}
\usepackage{xcolor}

\title{Memorization and Compositional Creativity in Small Language Models:\\
A Controlled Canary Study Across Transformer Architectures}

\author{\IEEEauthorblockN{Aarav Khanal}}

\begin{document}
\maketitle

\begin{abstract}
This working manuscript is not yet ready for submission. The exploratory study trains approximately
14-million-parameter decoder-only language models from scratch while manipulating the frequency of planted
narrative canaries. It compares near-parameter-matched positional-encoding, attention, and normalization variants
and separates reference-set exposure from global argmax extraction. The full automated grid and matching-corpus
exact four-gram audit are complete (50 cells and 16,000 generations). Human V/T/N labels have not been collected,
and the separate contamination-free confirmatory study has not been run. The evidence is exploratory and does not
establish a memorization-versus-creativity trade-off.
\end{abstract}

\begin{IEEEkeywords}
language-model memorization, compositional generalization, text generation, canary exposure, reproducibility
\end{IEEEkeywords}

\section{Introduction}
Language models can both recall training sequences and productively recombine familiar elements. Whether repeated
exposure shifts a small model from composition toward reproduction, and whether architecture changes that
relationship, remains an empirical question. We study that question through controlled from-scratch training,
rather than by inferring exposure from an opaque web corpus. We make no assumption that memorization must harm
composition; recent work also motivates neutral and beneficial relationships.

The intended contributions are: (1) a controlled narrative-canary exposure manipulation; (2) comparisons across
near-parameter-matched Transformer architectures; (3) separate reporting of reference-set exposure, global
argmax extraction, corpus-level n-gram overlap, diversity, and human judgments; and (4) an auditable exploratory
record with configurations, result files, tests, and runtime/failure notes. The current Git artifact does not
include the source corpus, tokenizer assets, or model checkpoints, so clean-room retraining is not yet possible.

\section{Related Work}
The exposure method follows Carlini et al.~\cite{carlini2019secret}. RAVEN demonstrates that novelty requires
comparison with training data at multiple linguistic scales and manual analysis~\cite{mccoy2023raven}. Haviv
et al. characterize recall of memorized idioms~\cite{haviv2023idioms}. Story-based memorization constraints have
also been studied in larger language models~\cite{jiang2026stories}, while Zhou et al. show that memorization can
support generalization through composition~\cite{zhou2026longtail}. Consequently, our contribution is not the
broad claim that memorization and composition have never been studied together. It is the controlled conjunction
of exposure frequency, small-model training, and near-fixed-size architecture ablations.

\section{Methods}
\subsection{Models and Intervention}
The baseline is a decoder-only model with eight layers, width 384, six query heads, two key--value heads, rotary
position embeddings, RMSNorm, SwiGLU, a 4096-token byte-level BPE vocabulary, and context length 512. The four
conditions are baseline (14,162,304 parameters), learned absolute positions (14,358,912), multi-head attention
with seven layers (13,964,928), and LayerNorm (14,162,304). Each checkpoint is trained from random initialization
for 6,000 steps, or 49.152 million training tokens per checkpoint, on the shared TinyStories-derived corpus.

Narrative canaries are inserted at 1, 10, 50, 100, or 500 stories; a zero-insertion condition is used as a matched
control, giving six conditions per architecture. Insertion pools are nested across frequencies and disjoint across canaries. Three originally selected
canaries were later found to overlap the base corpus. They are retained for auditability but excluded from the
clean-canary sensitivity aggregation; a separate confirmatory study has not yet generated its proposed new bank.

\subsection{Scale Transfer}
A separate extension trains the corresponding large-preset models (48.100M baseline, 48.428M learned-position,
47.239M MHA with nine layers, and 48.100M LayerNorm) on the same frequency-specific research corpora. Every
condition uses 6,000 updates with 8,192 tokens per update, or 49.152M tokens, to match the 14M arm's token budget.
The complete 24-cell large arm uses a fixed learning rate of 5e-4, following the existing large-model recipe; the
14M exploratory arm used 8e-4, so the scale comparison remains partly confounded by schedule choice. Every 48M
cell is single-seed and reuses the exploratory canary bank. The existing 48M MY-LLM checkpoint trained on
different data and 245.76M tokens and therefore is not treated as a matched canary experiment. The initial
48M baseline/frequency-0 attempt failed during checkpoint serialization and recovered from its verified step-1,000
checkpoint; the failure and recovery are retained in the runtime audit.

\subsection{Outcomes}
Memorization is measured by rank-based exposure relative to a fixed alternative set, with lift from the matched
zero-insertion model as the primary quantity. Teacher-forced per-token argmax matching tests a stricter extraction
threshold. Novelty evaluation includes distinct-$n$, a self-BLEU proxy, exact maximum word-four-gram overlap
against every matching training corpus, and a condition-blinded V/T/N judgment: verbatim reproduction, template
copy, or genuine novel recombination. Two blank 1,280-item rater sheets are prepared; no human labels, agreement,
or adjudication results exist yet. The exact word-four-gram audit is a copying diagnostic, not a semantic originality test.

\subsection{Analysis}
The exploratory grid is reported as exploratory. The confirmatory protocol uses mixed frequencies within one
corpus per architecture and seed, at least three seeds, a larger exposure reference set, clustered bootstrap
intervals, hierarchical repeated-measures models, multiplicity correction, and preregistered exclusions. Full
details are versioned with the artifact.

\section{Results}
All 24 canonical 14M cells, two additional baseline training-seed controls, and all 24 48M scale-transfer cells
completed 6,000 updates. Exact maximum word-four-gram overlap was evaluated against the matching frequency-specific
training corpus for 50 cells and 16,000 generations; corpus fingerprints and story counts were checked. The
baseline true-training-seed sensitivity analysis at frequencies 1, 10, 100, and 500 found clean-canary mean-lift
changes of $-0.588$, $+0.564$, $-0.636$, and $-0.436$ exposure units, respectively. The largest observed
per-canary changes ranged from 1.862 to 3.044. With two seeds and one architecture, these are descriptive
sensitivity results, not a variance estimate.

Table~\ref{tab:clean14} reports clean-canary matched 14M summaries. Exposure lift varies across frequency and
architecture, including plateaus and reversals. Distinct-2 does not provide a consistent monotonic novelty-cost
pattern, but this does not establish that a cost is absent: distinct-2 is not the human recombination endpoint,
and most cells have one seed. In the audited baseline and learned-position subset at five nonzero frequencies,
teacher-forced argmax extraction was 0/120 canary-checkpoint pairs. This bounded negative result does not
generalize beyond that subset or decoding criterion.

\begin{table*}[t]
\centering
\caption{Clean-canary descriptive summaries for the canonical 14M grid. Each cell is exposure lift over the
architecture-matched zero-insertion model / distinct-2 proxy. Most architecture-frequency conditions are one
training seed. These are descriptive summaries, not inferential estimates of a creativity effect.}
\label{tab:clean14}
\resizebox{\textwidth}{!}{%
\begin{tabular}{lcccccc}
\toprule
Architecture & 0x & 1x & 10x & 50x & 100x & 500x \\
\midrule
Baseline & 0.000/0.793 & 0.278/0.803 & 1.379/0.827 & 1.979/0.790 & 2.579/0.808 & 2.262/0.797 \\
Learned positions & 0.000/0.785 & 0.061/0.777 & 1.887/0.784 & 2.151/0.791 & 2.351/0.799 & 2.351/0.805 \\
MHA & 0.000/0.772 & 0.370/0.792 & 1.365/0.808 & 1.965/0.825 & 1.765/0.805 & 2.229/0.804 \\
LayerNorm & 0.000/0.817 & -0.094/0.792 & 2.015/0.808 & 2.132/0.788 & 1.932/0.782 & 1.932/0.799 \\
\bottomrule
\end{tabular}}
\end{table*}

The corresponding 48M grid is also complete, but single-seed, uses a different learning rate, and is exploratory.
Its matched descriptive table and figure are in \texttt{results/tradeoff\_clean.csv} and
\texttt{results/tradeoff\_clean.png}; they do not identify a causal scale effect. Per-cell runtimes and the
recovered checkpoint-write failure are reported in \texttt{results/trainer\_runtime\_audit.md}. No human novelty
score is available, and the confirmatory protocol in \texttt{docs/CONFIRMATORY\_STUDY\_PROTOCOL.md} has not been run.

\section{Limitations and Ethics}
The study is small in model and corpus scale. Most factorial cells and every 48M cell are single-seed; two
additional baseline seeds do not replicate the grid. The 48M learning rate differs from 14M, the design contains
post-registration changes, and three canary identities occurred in the base corpus. Diversity and n-gram
statistics are not substitutes for semantic novelty. Human annotation has not started; external annotation must
wait for any required institutional determination and consent. The confirmatory design requires power analysis
and other start-gate checks before any training. The final paper must report funding, conflicts, data/code
availability, and every exploratory deviation. The current Git artifact omits ignored source data and checkpoints;
a clean-environment reproduction and permanent archival are still needed.

\section{Conclusion}
The completed runs provide an exploratory audit of controlled canary exposure, not a settled answer about narrative
recombination. Human labels, an adequately powered contamination-free confirmation, and a reproducible public
artifact are still required before making a stronger claim.

\section*{Acknowledgment and AI-use disclosure}
AI-assisted tools were used for repository auditing, code and test drafting, literature discovery, and manuscript
editing. The human author is responsible for the research decisions, execution, verification, interpretation,
citations, and final text. This statement must be adapted to the exact policy of the selected venue and the final
record of tool use.

\begin{thebibliography}{00}
\bibitem{carlini2019secret} N. Carlini, C. Liu, {\'U}. Erlingsson, J. Kos, and D. Song, ``The Secret Sharer:
Evaluating and testing unintended memorization in neural networks,'' in \emph{28th USENIX Security Symposium},
2019, pp. 267--284.
\bibitem{mccoy2023raven} R. T. McCoy, P. Smolensky, T. Linzen, J. Gao, and A. Celikyilmaz, ``How much do language
models copy from their training data? Evaluating linguistic novelty in text generation using RAVEN,''
\emph{Transactions of the Association for Computational Linguistics}, vol. 11, pp. 652--670, 2023.
\bibitem{haviv2023idioms} A. Haviv, I. Cohen, J. Gidron, R. Schuster, Y. Goldberg, and M. Geva, ``Understanding
Transformer memorization recall through idioms,'' in \emph{Proc. EACL}, 2023, pp. 248--264.
\bibitem{jiang2026stories} Y. Jiang and F. Ferraro, ``Beyond math: Stories as a testbed for
memorization-constrained reasoning in LLMs,'' in \emph{Proc. EACL}, 2026, pp. 5590--5607.
\bibitem{zhou2026longtail} M. Zhou, H. Ma, and R. Ge, ``Memorizing long-tail data can help generalization through
composition,'' in \emph{Proc. ICLR}, 2026.
\end{thebibliography}

\end{document}
